\documentclass{article}
\usepackage{amsmath,amssymb}
\begin{document}
\title{Affine numerators in a shared Boolean filter}
\maketitle

Let $S$ have size $c$, let all $w_i>0$, and suppose $d>\sum_{i\in S}w_i$.
Write
\[
 K_S(d)=\sum_{T\subseteq S}(-1)^{c-|T|}
       \frac1{d-\sum_{i\in T}w_i},\qquad
 F_S(d)=\sum_{T\subseteq S}(-1)^{c-|T|}
       \frac{a+\sum_{i\in T}v_i}{d-\sum_{i\in T}w_i}.
\]
The exact identity
\[
 F_S(d)=\left(a+\sum_{i\in S}v_i\right)K_S(d)
             +\sum_{i\in S}v_iK_{S\setminus\{i\}}(d)
\]
is kernel checked as
\texttt{jointCell\_affine\_filtered\_kernel\_eq} in
\texttt{JointCellAffineFilteredNumerators.lean}.
The same file proves positivity when $a>0$ and all $v_i\geq0$.
All original definitions are unchanged.

This gives $c+1$ scalar features from the numerator of one main reciprocal
before subtracting its proper-face values. It does not supply $c+1$
independently chosen coefficients per column: the row coefficients $v_i$
are shared across columns, while the denominator offset and the column
numerator may vary. Pairing with a reference head gives structural zero on
the column axis at twice the main-head cost. The existing shared monomial
correction construction can fit the proper faces at the one-hot columns;
preserving arbitrary multihot prior inputs still requires the two-stage
gluing order recorded in the tensor audit.

The Laplace density is
\[
 \prod_{i\in S}(e^{w_i u}-1)
 \left(a+\sum_{i\in S}\frac{v_i e^{w_i u}}{e^{w_i u}-1}\right).
\]
Consequently its bracket belongs to the span of $1$ and
$\coth(w_i u/2)$. A Chebyshev theorem for this bracket cannot simply be
assumed: the separate exact ten-function coth determinant certificate
already refutes the all-orders version of that claim. Nor does a density
counterexample automatically survive the Laplace transform.

\texttt{JointCellAffineFilterSearch.py} records exploratory Wronskians of
the actual reciprocal feature family. Ten random six-coordinate examples
and four random ten-coordinate examples, each tested at nine exterior
distances, showed no sign reversal. Nearly equal weights
$w_i=1+i/1000$ at orders $8,10,12$, tested with 450 decimal digits, also
showed no reversal at eleven exterior distances. These are finite
floating-point computations, not sign certificates, a zero theorem, or a
decoder. No family with exponentially many strict sign words has been
constructed. Even a modest violation of a dimension-sized zero bound
would not by itself give the required unbounded saving.

\end{document}
